\documentclass[11pt]{article}

%% Page geometry
\usepackage[a4paper, margin=1in]{geometry}

%% Core math + fonts
\usepackage{amsthm,amsmath,amsfonts,amssymb}
\usepackage{microtype}

%% Citations
\usepackage[authoryear, round]{natbib}

%% Hyperlinks
\usepackage[colorlinks, citecolor=blue, urlcolor=blue, linkcolor=blue]{hyperref}
\usepackage{xurl}

%% Graphics & figures
\usepackage{graphicx}
\graphicspath{{./images/}}
\usepackage{float}
\usepackage{subcaption}
\usepackage{booktabs}

%% Algorithms
\usepackage[linesnumbered, ruled, vlined]{algorithm2e}

%% TikZ for diagrams
\usepackage{tikz}
\usetikzlibrary{shapes,arrows,fit,calc,positioning}
\tikzset{box/.style={draw, diamond, thick, text centered, minimum height=0.5cm, minimum width=1cm}}
\tikzset{line/.style={draw, thick, -latex'}}

%% Misc
\usepackage{comment}
\usepackage{dsfont}

%% Theorem environments
\theoremstyle{plain}
\newtheorem{axiom}{Axiom}
\newtheorem{claim}[axiom]{Claim}
\newtheorem{theorem}{Theorem}[section]
\newtheorem{lemma}[theorem]{Lemma}

\theoremstyle{definition}
\newtheorem{definition}[theorem]{Definition}
\newtheorem*{example}{Example}
\newtheorem*{fact}{Fact}

\theoremstyle{remark}
\newtheorem{case}{Case}

%% Custom math commands
\newcommand{\E}{\mathbb{E}}
\newcommand{\1}{\mathds{1}}
\newcommand{\p}{\mathbb{P}}
\renewcommand{\H}{\mathcal{H}}
\newcommand{\T}{\mathcal{T}}
\newcommand{\ind}{\perp\!\!\!\perp}

%% File-specific additions
\usepackage{bm}
\usepackage{mathtools}

\newcommand{\Var}{\mathrm{Var}}
\newcommand{\R}{\mathbb{R}}
\newcommand{\given}{\,|\,}
\newcommand{\indep}{\perp\!\!\!\perp}

%% Residual-prior abbreviations.  The map between these math symbols and the
%% R argument values is given once in Table~\ref{tab:sim-prior-grid}.
\newcommand{\mGauss}{\ensuremath{\mathrm{G}}}
\newcommand{\mShDP}{\ensuremath{\mathrm{sDP}}}
\newcommand{\miDP}{\ensuremath{\mathrm{iDP}}}
\newcommand{\miDPs}{\ensuremath{\mathrm{iDP}_\sigma}}
\newcommand{\mHDP}{\ensuremath{\mathrm{HDP}}}
\newcommand{\mHDPs}{\ensuremath{\mathrm{HDP}_\sigma}}


\title{ADDITIONAL SIMULATIONS}

\author{
  Tijn Jacobs}

\date{\today}

\setlength{\parindent}{0pt}
\setlength{\parskip}{0.8em}


\begin{document}
\maketitle

This file contains some specific simulations for the Supplementary Materials.
They should answer specific questions, or investigate properties of specific parts of the model.



\section{Simulation on the nonparametric error distribution}
\label{sec:sim-error-dist}

\subsection{Goal}
\label{sec:sim-goal}

The methodology section introduces a centered hierarchical Dirichlet
process (HDP) prior on the residual distribution, with the option of a
source-specific scale parameter $\sigma_s$. This simulation has three
aims.

\begin{enumerate}
  \item \textbf{Demonstrate the practical benefit of the
    \mHDPs prior.} Under data-generating processes
    where the two sources share residual shape but differ in
    scale---and especially when one source is too small to fit a
    flexible nonparametric residual on its own---atom sharing should
    yield lower CATE estimation error, narrower credible intervals
    without sacrificing coverage, and a more accurate per-source
    residual distribution than each of the five alternative residual
    priors.
  \item \textbf{Map out where the various priors do and do not
    differ.} The DGPs are designed so that each successive scenario
    breaks an assumption that one of the simpler priors relies on:
    Gaussianity, equality of source scales, equality of source shapes,
    and finally rarity of a residual ``type'' in one source but not the
    other. We expect a clean staircase of relative performance, with
    \mHDPs either tied for best or strictly best
    in every cell.
  \item \textbf{Probe the identifiability of the hierarchical
    parameters.} The atoms $\theta_k^*$ and the per-source centering
    shifts $\mu_s$ are not separately identified---only the centered
    locations $\theta_k^* - \mu_s$ are. We track posterior summaries
    of $|\mu_0 - \mu_1|$ and of the maximum unconstrained atom
    magnitude $\max_k |\theta_k^*|$ to confirm that this tension
    affects component-level interpretation without degrading the
    CATE estimates that the model targets.
\end{enumerate}

\subsection{Estimands and metrics}
\label{sec:sim-metrics}

For each replicate we record the following summaries.

\paragraph{CATE estimation.}
Posterior-mean root mean squared error
$\mathrm{RMSE}_\tau = \big( n^{-1} \sum_i (\hat\tau(X_i) - \tau(X_i))^2 \big)^{1/2}$,
bias $\mathrm{Bias}_\tau = n^{-1} \sum_i (\hat\tau(X_i) - \tau(X_i))$,
nominal-95\% credible-interval coverage averaged over $i$, and mean
credible-interval width. CATE estimates are evaluated on the training
sample with $\tau(x) = 1.5 + 2 x_1$.

\paragraph{Per-source residual scale recovery.}
We report two summaries.
\begin{itemize}
  \item \emph{Within-component scale.} The posterior mean of the
    within-component standard deviation $\sigma$ (or $\sigma_s$ in the
    scale modes). This is directly comparable to the truth only when
    the truth is Gaussian; under mixture truths the within-component
    $\sigma$ is intentionally smaller than the total residual sd, so
    cross-mode comparison on this metric is misleading.
  \item \emph{Total residual standard deviation.} The posterior mean of
    \[
      \hat\sigma^{\text{tot}}_s
      = \sqrt{\,\hat\sigma_s^2 + \textstyle\sum_k \hat\pi_{sk}\bigl(\hat\theta_{sk} - \hat\mu_s\bigr)^2\,}
    \]
    compared against the empirical standard deviation of the
    simulated source-specific residuals. This is the apples-to-apples
    metric across all six priors.
\end{itemize}

\paragraph{Per-source residual shape recovery.}
Two-sample Kolmogorov--Smirnov statistic between the empirical CDF of
the posterior-mean residuals $\{Y_i - \hat Y_i : S_i = s\}$ and the
true residuals $\{\varepsilon_i : S_i = s\}$ drawn during data
generation. Low values indicate that the prior is recovering the
source-specific residual law.

\paragraph{Mixture diagnostics.}
\begin{itemize}
  \item \emph{Effective number of components} per source,
    $K^{\text{eff}}_s = 1 / \sum_k \pi_{sk}^2$, averaged over MCMC
    iterations; the discrete analogue is the number of components
    with $\pi_{sk} > 0.01$.
  \item \emph{Atom sharing}: the number of components with both
    $\pi_{0k} > 0.01$ and $\pi_{1k} > 0.01$ (per iteration, averaged).
    Only meaningful for the source-stratified priors
    (\miDP, \miDPs, \mHDP, \mHDPs).
  \item \emph{Concentration parameters}: posterior means of the
    per-source concentration $M_s$ and, for the HDP priors, the
    top-level concentration $\gamma$.
\end{itemize}

\paragraph{HDP identifiability probes.}
For the HDP priors we additionally record the posterior means of
$|\mu_0 - \mu_1|$ and of $\max_k |\theta_k^*|$. Both quantities should
remain on the same scale as the true mean shift implied by the DGP and
the atom-prior scale $\sigma_\theta$, respectively; large values flag
the $\mu_s$--$\theta_k^*$ identifiability tension noted in
Section~\ref{sec:hdp-centering}.

\subsection{Compared priors}
\label{sec:sim-compared}

We compare six residual priors. They share a common skeleton and
differ along three axes:
\begin{enumerate}
  \item \textbf{Shape}: parametric Gaussian vs.\ nonparametric
    Dirichlet-process (DP) mixture of Gaussians;
  \item \textbf{Stratification}: residual distribution pooled across
    sources vs.\ stratified by source;
  \item \textbf{Borrowing}: when stratified, whether information is
    borrowed across sources via a hierarchical DP (HDP);
\end{enumerate}
with an orthogonal fourth axis for the residual scale:
\begin{enumerate}
  \setcounter{enumi}{3}
  \item \textbf{Scale}: a single residual scale $\sigma$ shared by
    both sources vs.\ source-specific scales $\sigma_0, \sigma_1$.
\end{enumerate}
We first set up the common building block, then state each prior as a
single small modification of the previous one.

\paragraph{Common building block: centered DP mixture of Gaussians.}
A Dirichlet-process mixture of Gaussians \citep{escobar1995bayesian}
models a residual as a mixture
\[
  f(w \given G, \sigma) \;=\; \int \frac{1}{\sigma}\, \phi\!\left(\frac{w - \theta}{\sigma}\right) dG(\theta),
\qquad
  G \sim \mathrm{DP}(\alpha, H),
\]
with $\phi$ the standard normal density and base measure
$H = N(0, \sigma_\theta^2)$. For computation we use the
Ishwaran--James truncation \citep{ishwaran2001gibbs} at $K$
components, so $G = \sum_{k=1}^K \pi_k\, \delta_{\theta_k}$ with
$\pi \sim \mathrm{Stick}_K(\alpha)$ and
$\theta_k \stackrel{\mathrm{iid}}{\sim} N(0, \sigma_\theta^2)$.
Mean-zero centering \citep{yang2010semiparametric} replaces the atoms
$\theta_k$ by $\theta_k - \mu$ with
$\mu = \sum_k \pi_k \theta_k$, ensuring
$\E[\varepsilon] = 0$ so that the structural-mean components
$m_0, \tau, c$ retain their interpretation.

\paragraph{1.\ \mGauss (baseline).}
A single Gaussian residual, shared across sources:
\[
  \varepsilon_i \given \sigma \;\sim\; N(0, \sigma^2).
\]
$\sigma$ has the usual scaled-inverse-$\chi^2$ prior.

\paragraph{2.\ \mShDP (one DP, pooled across sources).}
Replace the Gaussian by a single centered DP mixture, pooled across
both sources:
\[
  \varepsilon_i \given G, \sigma \;\sim\; f(\,\cdot\, \given G, \sigma)
  \quad\text{for all }i,
\]
with $G$ as in the common building block. Residual shape is allowed
to be non-Gaussian but is assumed identical across sources.

\paragraph{3.\ \miDP (per-source DP, independent).}
Fit a separate centered DP mixture per source, with no cross-source
borrowing. Two independent mixing measures $G_0, G_1$, common
$\sigma$:
\[
  \varepsilon_i \given S_i = s,\, G_s,\, \sigma \;\sim\;
   f(\,\cdot\, \given G_s, \sigma),
\quad
   G_s = \sum_{k=1}^K \pi_{sk}\, \delta_{\theta_{sk}}
   \text{ independently across } s.
\]
Sources can now differ in residual \emph{shape}.

\paragraph{4.\ \miDPs (per-source DP + per-source scale).}
Identical to \miDP but with source-specific scales:
\[
  \varepsilon_i \given S_i = s,\, G_s,\, \sigma_s \;\sim\;
   f(\,\cdot\, \given G_s, \sigma_s),
\]
with independent scaled-inverse-$\chi^2$ priors on each $\sigma_s^2$.

\paragraph{5.\ \mHDP (HDP, shared $\sigma$).}
The methodology section's prior. Replace the independent $G_0, G_1$
of \miDP by a hierarchical DP
\citep{teh2006hierarchical} that ties them through a top-level
discrete measure:
\[
  G^* \sim \mathrm{DP}(\gamma, H), \qquad
  G_s^* \sim \mathrm{DP}(M_s, G^*)
  \quad\text{for } s \in \{0, 1\}.
\]
Equivalently, after truncation,
$G^* = \sum_k \beta_k\, \delta_{\theta_k^*}$ has \emph{shared atoms}
$\theta_k^*$, and each per-source measure
$G_s^* = \sum_k \pi_{sk}\, \delta_{\theta_k^*}$ shares those atoms
but re-weights them by source-specific $\pi_{sk}$. After centering,
\[
  \varepsilon_i \given S_i = s,\, G_s^*,\, \sigma \;\sim\;
   \sum_{k=1}^K \pi_{sk}\,
   \phi(\,\cdot\,;\, \theta_k^* - \mu_s,\, \sigma),
\qquad
   \mu_s = \sum_k \pi_{sk}\, \theta_k^*.
\]
The shared atoms let evidence about the residual \emph{shape} flow
from the larger source to the smaller; the per-source weights and
the centering shifts $\mu_s$ retain source-specific flexibility.

\paragraph{6.\ \mHDPs (HDP + per-source scale).}
Layer the per-source scale of \miDPs onto the
HDP of \mHDP: shared atoms $\theta_k^*$, per-source
weights $\pi_{sk}$ and centering shifts $\mu_s$, \emph{and}
per-source scales $\sigma_s$. This is the prior the methodology
section adopts.

\paragraph{At-a-glance summary and notation key.}
Table~\ref{tab:sim-prior-grid} arranges the six priors along the
shape / stratification / borrowing axis (rows) and the
scale-sharing axis (columns). The math symbols in the second and
fourth columns are the abbreviations used throughout the remainder
of the document; the third and fifth columns give the corresponding
value of the \texttt{error\_dist} argument of
\texttt{FusionForest}, for use when running the code.

\begin{table}[tb]
\centering
\caption{The six residual priors compared in this simulation,
arranged by residual-shape construction (rows) and scale
parameterisation (columns). For each prior we give the math symbol
(used throughout this document) and the R argument value (used in
the \texttt{FusionForest()} call).}
\label{tab:sim-prior-grid}
\begin{tabular}{lcccc}
\toprule
& \multicolumn{2}{c}{Shared $\sigma$} & \multicolumn{2}{c}{Per-source $\sigma_s$} \\
\cmidrule(lr){2-3} \cmidrule(lr){4-5}
Residual shape & symbol & R argument & symbol & R argument \\
\midrule
Gaussian                                     & \mGauss & \texttt{gaussian}    & ---     & --- \\
Pooled DP mixture                            & \mShDP  & \texttt{shared\_dp}  & ---     & --- \\
Per-source DP mixtures, independent          & \miDP   & \texttt{source\_dp}  & \miDPs  & \texttt{source\_dp\_scale} \\
Per-source DP mixtures with HDP atom-sharing & \mHDP   & \texttt{source\_hdp} & \mHDPs  & \texttt{source\_hdp\_scale} \\
\bottomrule
\end{tabular}
\end{table}

\subsection{Data-generating process}
\label{sec:sim-dgp}

We generate observational and randomised data with covariates
$X \in [0,1]^5$, treatment $A \in \{0,1\}$, source indicator
$S \in \{0,1\}$ (1 = RCT, 0 = RWD), and continuous response
$Y = m_0(X) + A \tau(X) + (1 - S) U + \varepsilon$, where:
\begin{align*}
  m_0(X)    &= 2 X_1 - X_3^2, \\
  \tau(X)   &= 1.5 + 2 X_1, \\
  U         &\sim N(0, 1) \text{ in the RWD only (unmeasured confounder)}, \\
  A \given S = 1 &\sim \mathrm{Bernoulli}(0.5), \\
  A \given S = 0 &\sim \mathrm{Bernoulli}(\sigma_{\text{logit}}(0.8\, U)).
\end{align*}
The residual $\varepsilon$ depends on the source $S$ and the
scenario.

\paragraph{Scenario A: SHARED.}
$\varepsilon \given S = s \sim N(0, 1)$ for both sources. Null case;
all priors should perform similarly.

\paragraph{Scenario B: SCALE.}
$\varepsilon \given S = 1 \sim N(0, 0.5^2)$,
$\varepsilon \given S = 0 \sim N(0, 1.5^2)$. Sources differ in spread
only.

\paragraph{Scenario C: SHAPE\,+\,SCALE.}
For each source $s$, $\varepsilon$ is a two-component Gaussian
location mixture on atoms $\{-1.5, +1.0\}$ with per-source weights
$\pi_{\text{rct}} = (0.20, 0.80)$, $\pi_{\text{rwd}} = (0.75, 0.25)$,
recentered to mean zero per source, plus a residual Gaussian with
$\sigma_{\text{rct}} = 0.5$, $\sigma_{\text{rwd}} = 1.5$. Sources
share the atom \emph{support} but differ in both weights and scale.

\paragraph{Scenario D: TAIL SHARING.}
Atoms $\{-3.0, +0.5\}$, per-source weights
$\pi_{\text{rct}} = (0.05, 0.95)$, $\pi_{\text{rwd}} = (0.40, 0.60)$,
recentered to mean zero per source, with $\sigma_{\text{rct}} = 0.4$
and $\sigma_{\text{rwd}} = 0.8$. The RCT alone barely sees the
left-tail atom---in expectation $0.05\, n_{\text{rct}}$ observations.
This is the regime where HDP atom-borrowing should pay off, because the
RWD informs the RCT about the location of a residual feature the RCT
would otherwise be too small to identify.

\subsection{Sample-size scan}
\label{sec:sim-sample-size}

Each scenario is run at three RCT sample sizes,
$n_{\text{rct}} \in \{50, 150, 300\}$, with $n_{\text{rwd}} = 300$
fixed throughout. We expect the relative advantage of
\mHDPs over \miDPs (its
closest non-borrowing competitor) to grow as $n_{\text{rct}}$ shrinks,
especially in scenarios C and D.

\subsection{MCMC and reproducibility}
\label{sec:sim-mcmc}

All forests use $M = 200$ trees for the prognostic component, $M = 100$
for the treatment effect, and $M = 200$ for the deconfounding
component, with the default priors of \texttt{FusionForest}. The DP
truncation level is $K = 50$ and the atom-prior scale is
$\sigma_\theta = 0.5$ on the standardised response scale (defaults).
Each chain runs $5000$ burn-in iterations followed by $5000$ post-burn-in
draws. Per-cell results are averaged over $M$ replications, with $M$
matched to the number of available cores so that each replication runs
on a dedicated core; each replication fits all six priors serially. We
target $M \geq 100$ replications per cell.

The full simulation script is at
\texttt{examples/simulation\_error\_dist\_hpc.R}; the entire posterior
state needed to compute every metric in
Section~\ref{sec:sim-metrics} is stored in a single tidy data frame and
saved to \texttt{simulation\_error\_dist\_hpc\_output.rds}.

\subsection{Hypotheses and reporting plan}
\label{sec:sim-hypotheses}

We will report results as four blocks of tables and figures, mirroring
the four metric families.

\paragraph{H1 (overall CATE).}
\mHDPs matches the best of the alternatives on
$\mathrm{RMSE}_\tau$ in scenario~A (no truth-misspecification regime),
strictly outperforms in scenarios~B--D, and never loses coverage.

\paragraph{H2 (sample-size scan).}
In scenarios C and D, the gap between \mHDPs and
\miDPs grows as $n_{\text{rct}}$ shrinks; in
scenario D at $n_{\text{rct}} = 50$ the gap is largest and the
RCT-only residual sd of \miDPs is biased toward the
tighter Gaussian centre, while \mHDPs recovers the
tail via the RWD-informed shared atom.

\paragraph{H3 (mixture parsimony).}
$K^{\text{eff}}_s$ and the number of occupied components are smaller
for HDP priors than for the independent DPs in every cell, reflecting
the shrinkage induced by the top-level $\gamma$-prior.

\paragraph{H4 (identifiability check).}
$|\mu_0 - \mu_1|$ and $\max_k |\theta_k^*|$ remain bounded on the scale
implied by the DGP and the atom prior. Where they drift larger, this
will be documented but does not affect CATE conclusions, because the
posterior of $\theta_k^* - \mu_s$ remains tight.

\subsection{Results}
\label{sec:sim-results}

We ran $M = 191$ replications per cell on a 192-core node
(one core held in reserve). All tables below report the mean across
replications.

\subsubsection*{CATE estimation}

Table~\ref{tab:sim-rmse} reports the RMSE of $\hat\tau$ across the
twelve $(\text{scenario}, n_{\text{rct}})$ cells. The headline result
is that the scale-aware priors (\miDPs and
\mHDPs) match each other on CATE in every cell
and dominate the equal-$\sigma$ priors as soon as the truth has
unequal source spreads (scenarios B--D), with relative RMSE
reductions of $25\text{--}40\%$. The advantage of HDP atom-sharing
over independent per-source DPs (\miDPs versus
\mHDPs) is essentially zero on CATE, but neither
is HDP penalised: it matches the best in every cell. Coverage of the
nominal 95\% credible intervals is high and stable across modes
($\geq 0.98$ throughout); intervals are correspondingly narrower for
the scale-aware modes than for \mGauss.

\begin{table}[tb]
\centering
\small
\caption{CATE RMSE by scenario, RCT sample size, and residual prior.
Lowest value per row in bold. See
Table~\ref{tab:sim-prior-grid} for the prior abbreviations.}
\label{tab:sim-rmse}
\begin{tabular}{llcccccc}
\toprule
Scen. & $n_{\text{rct}}$ & G & sDP & iDP & iDP$_\sigma$ & HDP & HDP$_\sigma$ \\
\midrule
A & 50  & 0.471 & 0.473 & 0.470 & 0.476 & 0.469 & \textbf{0.466} \\
A & 150 & 0.356 & 0.357 & 0.355 & 0.354 & 0.355 & \textbf{0.354} \\
A & 300 & \textbf{0.290} & \textbf{0.290} & 0.291 & 0.291 & 0.291 & 0.292 \\
\midrule
B & 50  & 0.495 & 0.494 & 0.412 & 0.388 & 0.409 & \textbf{0.383} \\
B & 150 & 0.338 & 0.317 & 0.261 & \textbf{0.256} & 0.266 & \textbf{0.256} \\
B & 300 & 0.253 & 0.224 & 0.210 & \textbf{0.210} & 0.214 & 0.211 \\
\midrule
C & 50  & 0.571 & 0.574 & 0.538 & 0.526 & 0.538 & \textbf{0.520} \\
C & 150 & 0.413 & 0.410 & 0.380 & \textbf{0.329} & 0.388 & 0.353 \\
C & 300 & 0.332 & 0.272 & 0.230 & \textbf{0.223} & 0.269 & 0.224 \\
\midrule
D & 50  & 0.561 & 0.568 & 0.426 & 0.445 & \textbf{0.422} & 0.431 \\
D & 150 & 0.393 & 0.350 & 0.238 & \textbf{0.228} & 0.277 & 0.240 \\
D & 300 & 0.308 & 0.220 & 0.186 & \textbf{0.184} & 0.200 & 0.186 \\
\bottomrule
\end{tabular}
\end{table}

\subsubsection*{Per-source residual-distribution recovery}

Table~\ref{tab:sim-totalsd} reports the apples-to-apples total
residual sd error at $n_{\text{rct}} = 150$. The equal-$\sigma$ priors
(\mGauss, \mShDP) systematically
under-estimate the wider RWD source by $0.5\text{--}1.2$ across
B--D---a substantively large miscalibration. The two scale-aware
priors recover both per-source residual sds within roughly $0.2$.
Kolmogorov--Smirnov distances on the per-source residual CDFs (not
shown) cluster tightly across modes
($\mathrm{KS}_{\text{rwd}} \approx 0.05\text{--}0.07$,
$\mathrm{KS}_{\text{rct}} \approx 0.07\text{--}0.13$) with the
scale-aware priors typically a touch lower on the RCT side, but the
gap is small.

\begin{table}[tb]
\centering
\small
\caption{Total residual-sd error
$\hat\sigma^{\text{tot}}_s - \mathrm{sd}(\varepsilon_s^{\text{true}})$
at $n_{\text{rct}} = 150$. Negative values: under-estimating the
residual spread. Prior abbreviations as in
Table~\ref{tab:sim-prior-grid}.}
\label{tab:sim-totalsd}
\begin{tabular}{lcccc cccc}
\toprule
& \multicolumn{2}{c}{Scen.\ A} & \multicolumn{2}{c}{Scen.\ B}
& \multicolumn{2}{c}{Scen.\ C} & \multicolumn{2}{c}{Scen.\ D} \\
\cmidrule(lr){2-3} \cmidrule(lr){4-5} \cmidrule(lr){6-7} \cmidrule(lr){8-9}
Mode & RWD & RCT & RWD & RCT & RWD & RCT & RWD & RCT \\
\midrule
G            & $-0.45$ & $-0.46$ & $-0.89$ & $\hphantom{-}0.11$ & $-1.15$ & $-0.42$ & $-1.20$ & $-0.16$ \\
sDP          & $-0.20$ & $-0.21$ & $-0.47$ & $\hphantom{-}0.53$ & $-0.72$ & $\hphantom{-}0.01$ & $-0.56$ & $\hphantom{-}0.48$ \\
iDP          & $-0.06$ & $-0.30$ & $-0.03$ & $-0.16$ & $-0.27$ & $-0.40$ & $-0.04$ & $-0.09$ \\
iDP$_\sigma$ & $\hphantom{-}0.25$ & $\hphantom{-}0.05$ & $\hphantom{-}0.17$ & $\hphantom{-}0.11$ & $\hphantom{-}0.13$ & $\hphantom{-}0.06$ & $\hphantom{-}0.10$ & $\hphantom{-}0.09$ \\
HDP          & $-0.25$ & $-0.34$ & $-0.02$ & $-0.02$ & $-0.37$ & $-0.40$ & $-0.04$ & $-0.02$ \\
HDP$_\sigma$ & $\hphantom{-}0.25$ & $\hphantom{-}0.07$ & $\hphantom{-}0.18$ & $\hphantom{-}0.18$ & $\hphantom{-}0.14$ & $\hphantom{-}0.11$ & $\hphantom{-}0.15$ & $\hphantom{-}0.20$ \\
\bottomrule
\end{tabular}
\end{table}

\subsubsection*{Mixture parsimony}

This is the clearest qualitative win for the HDP priors.
Table~\ref{tab:sim-keff} shows the effective number of mixture
components per source. HDP atom sharing reduces
$K^{\text{eff}}_s$ by a factor of three to five relative to the
independent DPs, and \mHDPs is the sparsest
overall (typically 2--3 effective components per source). The number
of jointly-occupied components $n_{\text{shared}}$ (not shown) is
also substantially smaller under HDP, reflecting genuine sharing
rather than coincidence.

\begin{table}[tb]
\centering
\small
\caption{Effective number of mixture components per source,
$K^{\text{eff}}_s = 1 / \sum_k \pi_{sk}^2$, at $n_{\text{rct}} = 150$.
HDP atom-sharing yields a 3--5$\times$ reduction.}
\label{tab:sim-keff}
\begin{tabular}{lcccc cccc}
\toprule
& \multicolumn{2}{c}{Scen.\ A} & \multicolumn{2}{c}{Scen.\ B}
& \multicolumn{2}{c}{Scen.\ C} & \multicolumn{2}{c}{Scen.\ D} \\
\cmidrule(lr){2-3} \cmidrule(lr){4-5} \cmidrule(lr){6-7} \cmidrule(lr){8-9}
Mode & RWD & RCT & RWD & RCT & RWD & RCT & RWD & RCT \\
\midrule
iDP          & 13.1 & 7.0 & 17.0 & 1.5 & 15.0 & 3.8 & 11.0 & 1.7 \\
iDP$_\sigma$ & 10.5 & 6.9 & 12.2 & 1.5 & 12.1 & 3.8 &  5.7 & 1.7 \\
HDP          &  3.8 & 2.2 &  8.1 & 1.0 &  5.9 & 1.3 &  4.7 & 1.1 \\
HDP$_\sigma$ &  2.4 & 2.1 &  2.6 & 1.1 &  2.7 & 1.7 &  2.2 & 1.2 \\
\bottomrule
\end{tabular}
\end{table}

\subsubsection*{HDP identifiability check}

Table~\ref{tab:sim-identifiability} reports the posterior means of
$|\mu_0 - \mu_1|$ and $\max_k |\theta_k^*|$ across all twelve cells
for the two HDP priors. Both quantities drift to magnitudes of $10$
and beyond---far above the prior atom scale
$\sigma_\theta = 0.5$ on the standardised response scale and the
true DGP-implied mean shifts (zero in A, of order one in C, of order
1.5 in D). This is the $\mu_s$\,/\,$\theta_k^*$ non-identifiability
in action: the likelihood is invariant under
$\theta_k^* \mapsto \theta_k^* + c$,
$\mu_s \mapsto \mu_s + c$, and the sampler drifts along that
manifold. Crucially, this drift does \emph{not} compromise the
quantities that matter for inference: CATE RMSE (Table~\ref{tab:sim-rmse})
is identical to that of \miDPs, total residual
sds (Table~\ref{tab:sim-totalsd}) are recovered correctly, and the
mixture diagnostics in Table~\ref{tab:sim-keff} are well behaved.

\begin{table}[tb]
\centering
\small
\caption{HDP identifiability diagnostics: posterior means of
$|\mu_0 - \mu_1|$ and $\max_k |\theta_k^*|$, on the original response
scale, for the two HDP priors. Both drift well above what the DGP and
the atom-prior scale ($\sigma_\theta = 0.5$ standardised) would
imply; CATE inference is unaffected.}
\label{tab:sim-identifiability}
\begin{tabular}{llcccc}
\toprule
& & \multicolumn{2}{c}{HDP} & \multicolumn{2}{c}{HDP$_\sigma$} \\
\cmidrule(lr){3-4} \cmidrule(lr){5-6}
Scen. & $n_{\text{rct}}$ & $|\mu_0 - \mu_1|$ & $\max_k|\theta_k^*|$
                       & $|\mu_0 - \mu_1|$ & $\max_k|\theta_k^*|$ \\
\midrule
A & 50  & 12.3 &  3.8 & 11.4 & 4.0 \\
A & 150 & 14.6 &  5.6 & 13.4 & 5.9 \\
A & 300 & 15.6 &  7.3 & 14.5 & 7.2 \\
B & 50  & 13.9 &  3.7 & 10.2 & 3.9 \\
B & 150 & 16.0 &  7.6 & 14.3 & 6.3 \\
B & 300 & 17.0 & 13.1 & 15.6 & 7.5 \\
C & 50  & 13.1 &  4.1 & 11.4 & 4.4 \\
C & 150 & 15.7 &  7.0 & 14.5 & 7.0 \\
C & 300 & 17.1 &  9.6 & 14.7 & 8.4 \\
D & 50  & 12.3 &  4.8 & 11.5 & 4.5 \\
D & 150 & 16.1 &  9.3 & 14.7 & 8.3 \\
D & 300 & 17.4 & 12.1 & 15.5 & 9.1 \\
\bottomrule
\end{tabular}
\end{table}

\subsubsection*{Revisiting the hypotheses}

\begin{itemize}
  \item \textbf{H1} (\mHDPs is the strict CATE
    winner): \emph{partially supported}. It is never beaten on CATE,
    matches the best in 11/12 cells, but ties rather than strictly
    dominates \miDPs. The two scale-aware priors
    are statistically indistinguishable on CATE in these scenarios.
  \item \textbf{H2} (HDP--DP gap grows as $n_{\text{rct}}$ shrinks):
    \emph{not supported in scenario D}. At $n_{\text{rct}} = 50$
    plain \mHDP (no per-source scale) is the best
    individual mode in scenario D, suggesting that with very few RCT
    observations the extra $\sigma_s$ parameter is harder to identify
    than the shared atoms it borrows. In scenarios B and C the gap
    between scale-aware modes and the rest does shrink with sample
    size, as expected.
  \item \textbf{H3} (mixture parsimony): \emph{fully supported}.
    HDP priors use 3--5$\times$ fewer effective components than
    independent DPs in every cell.
  \item \textbf{H4} (identifiability bounded): \emph{not supported,
    but with no inferential cost}. $|\mu_0 - \mu_1|$ and
    $\max_k |\theta_k^*|$ drift well above the prior scale, but CATE
    estimates and total residual sds are unaffected.
\end{itemize}

\paragraph{Headline takeaway.}
\mHDPs is the safe default among the six priors:
it never loses on CATE, it recovers per-source residual scales
correctly, and it does so with substantially fewer effective mixture
components than the alternatives. Its main practical cost is the
$\mu_s$\,/\,$\theta_k^*$ non-identifiability, which must be addressed
in reporting (see below) but does not affect inference on the
estimands of interest.

\paragraph{Comment on identifiability (for the discussion).}
The HDP construction identifies only the centered locations
$\theta_k^* - \mu_s$. Posterior summaries that condition on a specific
component label---in particular the unconstrained atom locations
$\theta_k^*$ and the absolute means $\mu_s$---are therefore not
intrinsically meaningful and can drift along the manifold
$\theta_k^* \mapsto \theta_k^* + c$, $\mu_s \mapsto \mu_s + c$ without
changing the likelihood. The CATE estimand $\tau$ is unaffected: it
depends on the conditional mean only through the centering constraint
$\E[\varepsilon \given S = s] = 0$. We recommend reporting and
plotting only quantities invariant to this drift---in particular the
per-source residual densities
$\hat f_s(w) = \sum_k \pi_{sk} \phi\!\left((w - \theta_k^* + \mu_s) / \sigma_s\right)/\sigma_s$
and the centered atoms $\theta_k^* - \mu_s$.

\section{Simulation on the prior specification}
\label{sec:sim-prior}

\begin{remark}[Historical note on the parametrisation in this section]
The Phase-2 grid below was conducted under an earlier variance-budget
parametrisation with per-forest shrinkage factors $(k_d, k_\tau, k_c)$ in
a \emph{tightness} convention (larger $k$ $\Rightarrow$ tighter prior).
The current methodology has since dropped this framework in favour of the
uniform Chipman calibration $\sigma_{h,f} = k_f/(2\sqrt{m_f})$ with
$k_f = 1$ across all forests; the present section is retained as the
record of how that default was arrived at.
The five Phase-2 configurations (Default, Efficiency, Robust, Equal-share,
$\tau$-deep) and the associated $(k_d, k_\tau, k_c)$ tuples should be read
as the experimental conditions of the simulation, not as the methodology
default.
\end{remark}

\subsection{Goal}
\label{sec:sim-prior-goal}

The methodology section introduces a variance-budget calibration of the four
BART forests $(m_0^{\mathrm{sh}}, d, \tau, c)$, in which a single budget $\omega^2$
is split among the forests through per-forest shrinkage scales $k_f$. This
simulation has two aims.
\begin{enumerate}
  \item \textbf{Pick a default $(k_d, k_\tau, k_c)$.} Among five candidate
    variance-budget allocations, identify the one that delivers the lowest
    CATE RMSE at nominal coverage across a grid of confounding scenarios.
  \item \textbf{Probe the bias--efficiency trade-off.} When the real-world
    data are confounded, a prior that is too permissive on $c$ leaks bias
    into $\tau$; a prior that is too tight forfeits efficiency. The grid
    spans both regimes so that the chosen default can be defended as the
    Pareto-frontier choice rather than as a per-scenario winner.
\end{enumerate}
Tree counts and depth-prior parameters are held at the methodology defaults
in this study; a one-at-a-time sensitivity sweep around the chosen default
is deferred to a follow-up.

\subsection{Estimands and metrics}
\label{sec:sim-prior-metrics}

The only estimand scored is the CATE on the log-time scale,
$\tau(x) = \E[\log T(1) - \log T(0) \mid X = x]$, evaluated
\emph{in-sample} on the combined training data
($n = n_1 + n_0$ points per replication: RCT plus RWD).  Bayesian shrinkage
from the BART prior plays the role that a hold-out set would play under
maximum-likelihood inference: the model is a regularised estimator of
$\tau$, and the in-sample posterior mean is exactly the quantity returned
to a downstream analyst, so in-sample evaluation is the most direct
operational test.  For each replicate we record the posterior mean and
the equal-tailed 95\% credible interval of $\tau$ at each training point
and compute four summaries, averaged over replications.
\begin{itemize}
  \item \emph{RMSE.}
    $\mathrm{RMSE}_\tau = \big( n^{-1} \sum_{i=1}^{n} (\hat\tau(X_i) - \tau(X_i))^2 \big)^{1/2}$.
  \item \emph{Pointwise 95\% coverage.}
    $n^{-1} \sum_{i=1}^{n} \1\{\tau(X_i) \in \mathrm{CI}_{95}(X_i)\}$;
    nominal target $0.95$.
  \item \emph{Mean credible-interval width.} An efficiency proxy.
  \item \emph{Bias.}
    $n^{-1} \sum_{i=1}^{n} (\hat\tau(X_i) - \tau(X_i))$ (signed);
    detects confounding leakage.
\end{itemize}
The confounding function $c$ is an internal nuisance and is not scored
directly.

\subsection{Compared prior configurations}
\label{sec:sim-prior-configs}

The leaf prior of each forest is $h_{j,a} \sim N(0, \omega_f^2)$ per leaf,
with $\omega_f = k_f^{\mathrm{code}} / \sqrt{m_f}$. The five configurations
differ in their setup-side $(k_d, k_\tau, k_c)$, in the \emph{tightness}
convention (larger $k$ $\Rightarrow$ tighter prior). The setup-side values
are mapped to the code-side leaf-prior scales by
\[
  s_0 = \bigl(1 + k_d^{-2} + k_\tau^{-2} + k_c^{-2}\bigr)^{-1/2},
  \qquad
  k_{\mathrm{sh}}^{\mathrm{code}} = s_0,
  \qquad
  k_f^{\mathrm{code}} = s_0 / k_f \quad (f \in \{d, \tau, c\}),
\]
which fixes the total prior variance
$\Var_\pi\bigl(m_0^{\mathrm{sh}}(X)\bigr) + \Var_\pi(d(X)) + \Var_\pi(\tau(X)) + \Var_\pi(c(X)) = 1$
on the standardised log-time scale. The depth columns in
Table~\ref{tab:sim-prior-configs} are $(\mathrm{base}, \mathrm{power})$,
i.e.\ the non-terminal probability $\mathrm{base}\,(1+d)^{-\mathrm{power}}$
of \citet{chipman1998bayesian}; the baseline and deviation forests use the
BART default $(0.95, 2)$ throughout.

\begin{table}[tb]
\centering
\small
\caption{The five prior configurations compared in this simulation. The
columns $k_d, k_\tau, k_c$ are in the tightness convention (larger
$\Rightarrow$ tighter); $k_{\mathrm{sh}} \equiv 1$ throughout. Trees are
listed in the order $(m_{\mathrm{sh}}, m_d, m_\tau, m_c)$. Depth columns
are $(\mathrm{base}, \mathrm{power})$.}
\label{tab:sim-prior-configs}
\begin{tabular}{lcccccl}
\toprule
Config       & $k_d$ & $k_\tau$ & $k_c$ & trees & depth $\tau$ / $c$ & role \\
\midrule
Default      & 2.5 & 1.7 & 4  & 200,\,50,\,200,\,200 & $(0.95, 3)$ / $(0.25, 3)$ & proposed          \\
Efficiency   & 4   & 1.7 & 10 & 200,\,50,\,200,\,200 & $(0.95, 3)$ / $(0.25, 3)$ & trust RWD         \\
Robust       & 2   & 1.7 & 2  & 200,\,50,\,200,\,200 & $(0.95, 3)$ / $(0.25, 3)$ & discount RWD      \\
Equal-share  & 1   & 1   & 1  & 200,\,50,\,200,\,200 & $(0.95, 3)$ / $(0.25, 3)$ & agnostic baseline \\
$\tau$-deep  & 2.5 & 1.7 & 4  & 200,\,50,\,200,\,200 & $(0.95, 2)$ / $(0.25, 3)$ & richer HTE        \\
\bottomrule
\end{tabular}
\end{table}

The Default configuration encodes a BCF-style ordering with the largest
variance on the baseline, a moderate budget on $\tau$, a tighter prior on
the RWD deviation $d$, and the tightest prior on the confounding function
$c$. The Default depth on $\tau$ is $(0.95, 3)$ to enforce shallow trees
\citep{hahn2020bayesian}; on $c$ the base parameter is dropped to $0.25$ to
enforce near-constant $c$ trees, motivated by the structural argument that
$c \equiv 0$ corresponds to no unmeasured confounding in the RWD.
Efficiency tightens the priors on $d$ and $c$ to lean on the RWD; Robust
loosens them; Equal-share is the agnostic 25\% split per forest and is
included to demonstrate that the BCF-style ordering is empirically
justified; $\tau$-deep relaxes the $\tau$ depth penalty to the BART default
to allow richer heterogeneity.

\subsection{Data-generating process}
\label{sec:sim-prior-dgp}

We generate observational and randomised data with covariates
$X \in [0,1]^{10}$, treatment $A \in \{0,1\}$, source indicator
$S \in \{0,1\}$ (1 = RCT, 0 = RWD), and continuous log-survival response
\[
  \log T \;=\; f_0(X, U) \;+\; A\, \tau^*(X, U) \;+\; \sigma\, \varepsilon,
  \qquad \varepsilon \sim N(0, 1).
\]
Five covariates $X_1, \ldots, X_5$ are active and five $X_6, \ldots, X_{10}$
are noise; all are drawn iid $U(0, 1)$ in both sources, with no
covariate-shift between RCT and RWD. Let $Z_j = X_j - 0.5$. The structural
components are
\begin{align*}
  f_0(X, U)    &= 0.7\, \sin(\pi Z_1) + Z_2 Z_3 + \bigl(Z_4^2 - \tfrac{1}{12}\bigr) + \beta_U\, U, \\
  \tau^*(X, U) &= \tau(X) + \gamma_U\, U, \\
  \tau(X)      &= \tau_0 + \tau_1\, Z_1 + \tau_2\, \1\{X_2 > 0.5\},
\end{align*}
with $U \sim N(0, 1)$ independent of $X$ and identically distributed across
sources. The unit-level effect $\tau^*$ depends on $U$ through $\gamma_U$;
since $U \indep X$, the CATE $\tau(x) = \E[\tau^*(X, U) \mid X = x]$ is
exactly the function $\tau(X)$ above. Two CATE regimes are studied: the
\emph{homogeneous} regime $(\tau_0, \tau_1, \tau_2) = (0.5, 0, 0)$ and the
\emph{heterogeneous} regime $(0.5, 0.4, 0.3)$. The noise scale is
$\sigma = 1$, giving a signal-to-noise ratio of approximately one on the
standardised log-time scale.

Treatment in the RCT is randomised, $A \mid S = 1 \sim \mathrm{Bernoulli}(0.5)$;
in the RWD treatment depends on observed and unmeasured covariates,
\[
  A \mid S = 0 \;\sim\; \mathrm{Bernoulli}\bigl(\mathrm{expit}(\alpha_0 + \alpha_X(Z_1 + Z_2) + \alpha_U\, U)\bigr),
\]
with $\alpha_X = 1$ and $\alpha_0$ tuned per scenario so that the marginal
treatment prevalence $\Pr(A = 1 \mid S = 0) \approx 0.5$. The parameter
$\alpha_U$ governs the strength of unmeasured confounding through the
$U$--$A$ association.

\paragraph{Confounding scenarios.}
Four scenarios span the $(\alpha_U, \beta_U, \gamma_U)$ space:
\begin{itemize}
  \item \textbf{S0 (no confounding):} $(0, 1, 0)$. $U \indep A \mid X, S = 0$,
    so the implied $c(x) \equiv 0$ and the RWD is unconfounded.
  \item \textbf{S1 (mild):} $(0.5, 1, 0)$. Small prognostic-imbalance term in
    $c$.
  \item \textbf{S2 (strong):} $(1.5, 1.5, 0)$. Large prognostic-imbalance
    term.
  \item \textbf{S3 (effect-modification):} $(1, 1, 0.5)$. Moderate $c$ with
    an added effect-modification component
    $\gamma_U\, \E[U \mid x, A = 1, S = 0]$.
\end{itemize}

\paragraph{Censored and uncensored variants.}
The simulation is run in two variants. In the \emph{censored variant} an
administrative threshold $c^\star$ is applied to every observation, with
$c^\star$ tuned once at script start (via a large pre-simulation) so that
the expected censoring rate equals $35\%$ under the central cell (S1,
heterogeneous, largest sample size); the same $c^\star$ is then used across
all cells. In the \emph{uncensored variant} the threshold is removed
($\delta_i \equiv 1$) and $\log T$ is standardised to unit variance via the
sample standard deviation before fitting. The uncensored variant serves as
a clean reference that removes the wobble introduced by the preliminary
AFT-based scale estimate that the censored fit relies on.

\subsection{Sample sizes and grid}
\label{sec:sim-prior-grid}

The full primary grid crosses four confounding scenarios, two CATE regimes
(homogeneous, heterogeneous), and two sample-size pairs
$(n_1, n_0) \in \{(250, 1000),\,(100, 400)\}$, yielding $4 \times 2 \times 2 = 16$
cells. Within each cell the five prior configurations of
Table~\ref{tab:sim-prior-configs} are fitted on the same simulated dataset
(common random numbers across configurations) so that performance
differences are not inflated by dataset-to-dataset noise.

\subsection{MCMC and reproducibility}
\label{sec:sim-prior-mcmc}

All forests use the tree counts in Table~\ref{tab:sim-prior-configs}. The
residual is Gaussian (\texttt{error\_dist = "gaussian"}) and treatment is
centred (\texttt{treatment\_coding = "centered"}). Each fit runs a
\emph{single} MCMC chain of $2000$ warm-up plus $2000$ retained
iterations, with no thinning; standardisation of $\log T$ is handled
inside each replicate, on its own data. We run $R = 300$ replications per
cell. Replications are parallelised across cores; each replication fits
all five configurations serially on its assigned core.

The HPC scripts driving the two variants are
\texttt{simulations/prior/simulation\_censored\_hpc.R} and
\texttt{simulations/prior/simulation\_uncensored\_hpc.R}. Output is written
to \texttt{simulation\_censored\_hpc\_output.rds} and
\texttt{simulation\_uncensored\_hpc\_output.rds} respectively, in long
form: one row per (scenario, CATE regime, sample-size pair, configuration,
replication) cell, carrying the four CATE metrics.

\subsection{Hypotheses and reporting plan}
\label{sec:sim-prior-hypotheses}

\paragraph{H1 (Default sits on the Pareto frontier).}
Across the four confounding scenarios, the Default configuration achieves
the lowest average $\mathrm{RMSE}_\tau$ subject to maintaining at least
$0.90$ coverage at every cell.

\paragraph{H2 (Efficiency wins in S0, loses in S2/S3).}
In S0, where the truth is $c \equiv 0$, the Efficiency configuration
($k_c = 10$) wins on $\mathrm{RMSE}_\tau$ because the tight prior on $c$
maximises the effective sample size for $\tau$. In S2 and S3, Efficiency
shows substantial bias and under-coverage, because the tight
prior on $c$ prevents the model from absorbing the RWD confounding.

\paragraph{H3 (Robust is conservative).}
Robust matches the Default in S2 and S3 but gives a smaller efficiency
gain over the agnostic Equal-share baseline in S0 and S1.

\paragraph{H4 ($\tau$-deep helps under heterogeneity, hurts otherwise).}
In the heterogeneous CATE regime, $\tau$-deep matches or beats the Default
on $\mathrm{RMSE}_\tau$; in the homogeneous regime it is slightly worse,
because deeper $\tau$ trees overfit the constant truth.

\paragraph{H5 (censored vs uncensored).}
The ranking of configurations is preserved between the censored and
uncensored variants. Absolute $\mathrm{RMSE}_\tau$ is higher under
censoring, reflecting the information loss from $35\%$ administrative
censoring and the additional uncertainty in the preliminary AFT-based
scale estimate.

\subsection{Results}
\label{sec:sim-prior-results}

We report preliminary results from $R = 190$ replications per cell,
$N_{\mathrm{burn}} = N_{\mathrm{post}} = 2000$, on the \emph{uncensored}
variant. The full $R = 300$ runs and the censored variant will replace
these numbers in a later revision; the relative ordering of the
configurations is expected to be stable. All tables in this subsection
average over the $190$ replications.

The homogeneous and heterogeneous CATE regimes produce essentially
identical relative rankings of the five configurations. We therefore
present the heterogeneous CATE as the primary result and report the
homogeneous CATE for completeness in
Table~\ref{tab:sim-prior-rmse-homog}.

\subsubsection*{CATE RMSE}

Table~\ref{tab:sim-prior-rmse-het} reports $\mathrm{RMSE}_\tau$ across the
eight $(\text{scenario}, n_1, n_0)$ cells of the heterogeneous regime. The
ranking is consistent across all four scenarios.

\begin{table}[tb]
\centering
\small
\caption{CATE RMSE by scenario, sample size, and configuration, for the
heterogeneous CATE regime. Lowest per row in bold.}
\label{tab:sim-prior-rmse-het}
\begin{tabular}{llccccc}
\toprule
Scen. & $(n_1, n_0)$  & Default & Efficiency      & Robust          & Equal-share     & $\tau$-deep \\
\midrule
S0    & $(100, 400)$  & 0.257   & \textbf{0.246}  & 0.276           & 0.297           & 0.257       \\
S0    & $(250, 1000)$ & 0.233   & \textbf{0.225}  & 0.234           & 0.247           & 0.235       \\
\midrule
S1    & $(100, 400)$  & 0.278   & 0.363           & \textbf{0.270}  & 0.296           & 0.278       \\
S1    & $(250, 1000)$ & 0.241   & 0.323           & \textbf{0.232}  & 0.249           & 0.242       \\
\midrule
S2    & $(100, 400)$  & 0.572   & 1.038           & \textbf{0.357}  & 0.377           & 0.571       \\
S2    & $(250, 1000)$ & 0.401   & 0.864           & \textbf{0.298}  & 0.316           & 0.402       \\
\midrule
S3    & $(100, 400)$  & 0.413   & 0.680           & \textbf{0.310}  & 0.337           & 0.412       \\
S3    & $(250, 1000)$ & 0.318   & 0.585           & \textbf{0.269}  & 0.287           & 0.319       \\
\bottomrule
\end{tabular}
\end{table}

\textbf{Robust} is the lowest-RMSE configuration in seven of the eight
heterogeneous cells and is beaten only in S0 by \textbf{Efficiency} (and
even there only by 0.01--0.03). Across S1--S3 the gap to the runners-up
grows: Robust beats Default by $25\text{--}40\%$ on RMSE in S2 and
$25\text{--}30\%$ in S3. Equal-share is the second-best configuration in
the confounded scenarios but consistently $0.03$--$0.04$ above Robust on
RMSE. \textbf{Efficiency} is the across-the-board worst configuration as
soon as confounding is present: its RMSE in S2 is roughly three times
that of Robust at both sample sizes. \textbf{$\tau$-deep} is
indistinguishable from Default in every cell (differences are at the
third decimal place), which means the depth-prior change on $\tau$ does
not materially affect CATE estimation at this signal-to-noise level.

The homogeneous-regime RMSE in Table~\ref{tab:sim-prior-rmse-homog}
mirrors this ordering, with absolute values $0.02$--$0.03$ below the
heterogeneous regime everywhere.

\begin{table}[tb]
\centering
\small
\caption{CATE RMSE by scenario, sample size, and configuration, for the
homogeneous CATE regime. Lowest per row in bold.}
\label{tab:sim-prior-rmse-homog}
\begin{tabular}{llccccc}
\toprule
Scen. & $(n_1, n_0)$  & Default & Efficiency      & Robust          & Equal-share     & $\tau$-deep \\
\midrule
S0    & $(100, 400)$  & 0.224   & \textbf{0.215}  & 0.236           & 0.266           & 0.225       \\
S0    & $(250, 1000)$ & 0.214   & \textbf{0.205}  & 0.211           & 0.229           & 0.216       \\
\midrule
S1    & $(100, 400)$  & 0.250   & 0.344           & \textbf{0.236}  & 0.273           & 0.251       \\
S1    & $(250, 1000)$ & 0.223   & 0.313           & \textbf{0.204}  & 0.224           & 0.224       \\
\midrule
S2    & $(100, 400)$  & 0.583   & 1.038           & \textbf{0.328}  & 0.342           & 0.581       \\
S2    & $(250, 1000)$ & 0.407   & 0.885           & \textbf{0.284}  & 0.303           & 0.407       \\
\midrule
S3    & $(100, 400)$  & 0.414   & 0.690           & \textbf{0.277}  & 0.306           & 0.414       \\
S3    & $(250, 1000)$ & 0.308   & 0.597           & \textbf{0.240}  & 0.260           & 0.309       \\
\bottomrule
\end{tabular}
\end{table}

\subsubsection*{Bias}

Signed bias (Table~\ref{tab:sim-prior-bias-het}) makes the failure mode
of Efficiency explicit. Under no confounding (S0) every configuration is
near zero ($|\mathrm{Bias}|\leq 0.10$). Under confounding the bias of
Efficiency, Default and $\tau$-deep climbs monotonically with the
confounding strength: in S2, Efficiency carries a bias of approximately
$+1.0$ in both heterogeneity regimes (the RWD treated-vs-control
contrast is approximately one log-time unit larger than the true CATE).
Default and $\tau$-deep carry roughly half that bias. \textbf{Robust}
and \textbf{Equal-share} are the only configurations that keep the bias
small under confounding; Equal-share is the smallest-bias option in S2
and S3, with Robust a close second.

\begin{table}[tb]
\centering
\small
\caption{CATE bias (signed integrated bias of $\hat\tau$) by scenario,
sample size, and configuration, for the heterogeneous CATE regime.
Smallest $|\mathrm{Bias}|$ per row in bold.}
\label{tab:sim-prior-bias-het}
\begin{tabular}{llccccc}
\toprule
Scen. & $(n_1, n_0)$  & Default & Efficiency       & Robust            & Equal-share        & $\tau$-deep \\
\midrule
S0    & $(100, 400)$  & $-0.057$ & \textbf{$-0.019$} & $-0.097$         & $-0.077$           & $-0.058$    \\
S0    & $(250, 1000)$ & $-0.036$ & \textbf{$-0.011$} & $-0.054$         & $-0.042$           & $-0.038$    \\
\midrule
S1    & $(100, 400)$  & $+0.105$ & $+0.272$         & \textbf{$-0.031$} & $-0.041$           & $+0.103$    \\
S1    & $(250, 1000)$ & $+0.067$ & $+0.237$         & \textbf{$-0.015$} & $-0.020$           & $+0.066$    \\
\midrule
S2    & $(100, 400)$  & $+0.498$ & $+1.000$         & $+0.133$          & \textbf{$+0.060$}  & $+0.495$    \\
S2    & $(250, 1000)$ & $+0.292$ & $+0.824$         & $+0.064$          & \textbf{$+0.023$}  & $+0.290$    \\
\midrule
S3    & $(100, 400)$  & $+0.303$ & $+0.626$         & $+0.053$          & \textbf{$+0.008$}  & $+0.300$    \\
S3    & $(250, 1000)$ & $+0.181$ & $+0.532$         & $+0.023$          & \textbf{$+0.000$}  & $+0.180$    \\
\bottomrule
\end{tabular}
\end{table}

The bias pattern explains the RMSE pattern: under confounding the squared
bias dominates the squared variance for every configuration except
Robust and Equal-share, so the configurations that fail to absorb the
RWD bias into $c$ pay for it in RMSE.

\subsubsection*{Coverage}

Table~\ref{tab:sim-prior-coverage-het} reports pointwise 95\%
credible-interval coverage. All configurations except Efficiency exceed
the nominal $0.95$ target everywhere (in fact they over-cover at
$\geq 0.94$ in every cell). The variance-budget calibration is therefore
generally conservative; intervals are wider than strictly needed but the
nominal level is not violated. \textbf{Efficiency} is the only
configuration that under-covers, and severely so: $0.51$--$0.54$ in S2
and $0.81$--$0.84$ in S3.

\begin{table}[tb]
\centering
\small
\caption{Pointwise 95\% credible-interval coverage by scenario, sample
size, and configuration, for the heterogeneous CATE regime. Nominal
target $0.95$; configurations under-covering are flagged in italics.}
\label{tab:sim-prior-coverage-het}
\begin{tabular}{llccccc}
\toprule
Scen. & $(n_1, n_0)$  & Default & Efficiency           & Robust  & Equal-share & $\tau$-deep \\
\midrule
S0    & $(100, 400)$  & 0.998   & 0.999                & 0.994   & 0.995       & 0.999       \\
S0    & $(250, 1000)$ & 0.996   & 0.997                & 0.994   & 0.996       & 0.997       \\
\midrule
S1    & $(100, 400)$  & 0.997   & 0.983                & 0.998   & 0.997       & 0.997       \\
S1    & $(250, 1000)$ & 0.996   & 0.976                & 0.996   & 0.996       & 0.997       \\
\midrule
S2    & $(100, 400)$  & 0.945   & \emph{0.512}         & 0.995   & 0.996       & 0.953       \\
S2    & $(250, 1000)$ & 0.966   & \emph{0.535}         & 0.991   & 0.994       & 0.971       \\
\midrule
S3    & $(100, 400)$  & 0.982   & \emph{0.835}         & 0.994   & 0.994       & 0.984       \\
S3    & $(250, 1000)$ & 0.986   & \emph{0.821}         & 0.995   & 0.996       & 0.989       \\
\bottomrule
\end{tabular}
\end{table}

Mean credible-interval widths (not tabulated) lie between $1.35$ and
$2.25$ across all cells and configurations and differ by less than
$10\%$ between any two configurations within a cell. Efficiency does
\emph{not} buy meaningful interval-width reduction over Robust, even
though it shrinks $c$ harder; the widths are nearly the same, so the
extra coverage loss in S2/S3 comes from a shift in the posterior centre
rather than from artificially tight intervals.

\subsubsection*{Revisiting the hypotheses}

\begin{itemize}
  \item \textbf{H1} (Default is Pareto-optimal): \emph{not supported}.
    Default is dominated by Robust in S1--S3 on both RMSE and bias and
    is essentially tied with Robust in S0. The choice of $k_c = 4$ in
    Default over-tightens the confounding prior; the data prefer the
    looser $k_c = 2$ of Robust.
  \item \textbf{H2} (Efficiency wins S0, loses S2/S3):
    \emph{fully supported}. Efficiency has the lowest RMSE in S0 by a
    small margin and the largest RMSE and largest bias in S1, S2 and S3
    with a coverage of only $0.51$--$0.54$ in S2.
  \item \textbf{H3} (Robust is the conservative choice):
    \emph{partially supported}. Robust is conservative in coverage
    ($\geq 0.99$ everywhere) but is also the lowest-RMSE configuration
    in seven of eight heterogeneous cells. ``Conservative'' over-states
    its cost.
  \item \textbf{H4} ($\tau$-deep helps under heterogeneity):
    \emph{not supported}. $\tau$-deep is indistinguishable from
    Default at three decimal places in every cell, in both CATE
    regimes. The depth-prior change on $\tau$ has no practical effect
    at this signal-to-noise level.
  \item \textbf{H5} (censored vs uncensored ranking preserved):
    \emph{to be checked} once the censored run completes; the ranking
    above is from the uncensored run only.
\end{itemize}

\paragraph{Headline takeaway (Phase 2).}
On this $R = 190$ preliminary run, \textbf{Robust} is the across-the-board
winner among the five Phase-2 configurations: lowest CATE RMSE in S1--S3,
near-best in S0, no coverage failures anywhere, and the smallest bias
outside of Equal-share (which is itself competitive).
The ``Default'' configuration $(k_d, k_\tau, k_c) = (2.5, 1.7, 4)$ is too
tight on $c$; the data prefer the looser $k_c = 2$ of Robust.
Within the variance-budget framework this would motivate adopting Robust
$(k_d, k_\tau, k_c) = (2, 1.7, 2)$ as the new default.

\paragraph{Caveat on coverage and the follow-up sweeps.}
All non-Efficiency configurations over-cover ($\geq 0.94$ everywhere).
This indicates the variance-budget calibration is too generous: the
budget $\omega^2 = 1$ on standardised $\log T$ allocates more prior
variance to the structural component than is coherent with the
inverse-$\chi^2$ residual prior.
Phases~2b (v2, $\omega \in \{0.95, 1\}$), 2d (v4, $\omega \in \{0.5, 0.75, 1\}$),
and 2e (v5, parametrisation comparison) trace this thread and converge on
the uniform Chipman calibration $\sigma_{h,f} = k_f/(2\sqrt{m_f})$ with
$k_f = 1$ for every forest as the adopted methodology default.
The asymmetric variance-budget allocation diagnosed here gave a small
RMSE advantage under strong confounding but at the cost of an opaque
calibration; the uniform Chipman default is simpler, matches the BART/BCF
defaults of the implementation, and is internally coherent with the
residual prior. See the prior-specification subsection of
\textbf{METHODOLOGY.tex} for the final form.

\bibliographystyle{apalike}
\bibliography{../general/references}
\end{document}
